% from ega4-errata-addenda-list3-fr.tex
\oldpage[IV]{357}
Ligne 5 de la p.~191, ajouter : \emph{Si de plus les fibres
$f^{-1}(f(x))$ de $f$ sont équidimensionnelles ($0_{\mathrm{IV}}$,
14.1.3), on dit que $f$ est équidimensionnel.} Ligne 6 du bas de la
p.~194, remplacer « équidimensionnelles » par « localement
équidimensionnelles ».

\paragraph{(Err\textsubscript{IV}, 35).}
Dans le titre de (IV, 13.3), ligne 4 du bas de la p.~94, remplacer
« équidimensionnels » par « localement équidimensionnels ». Dans
(IV, 13.3.2), ligne 6 du bas de la p.~196, remplacer deux fois
« équidimensionnel » par « localement équidimensionnel ». Ligne 5 du bas
de la p.~196, ajouter : \emph{Si de plus les fibres $f^{-1}(f(x))$ de
$f$ sont équidimensionnelles ($0_{\mathrm{IV}}$, 14.1.3), on dit que
$f$ est équidimensionnel.} Dans (IV, 13.3.7), ligne 14 de la p.~198,
remplacer « équidimensionnel » par « localement équidimensionnel ».
Dans ({\hyperref[canon-index.ega4.c04.statement.13-3-9]{IV, 13.3.9}}), ligne 3 du
bas de la p.~198, remplacer « équidimensionalité » par
« équidimensionalité locale ». Lignes 10, 12, 20 et 24 de la p.~199,
remplacer « équidimensionnel » par « localement équidimensionnel ».

\paragraph{(Err\textsubscript{IV}, 36).}
Dans l'introduction du \S14, ligne 6 du bas de la p.~199, remplacer
« équidimensionnel » par « localement équidimensionnel ».

\paragraph{(Err\textsubscript{IV}, 37).}
Dans (IV, 14.4.1), ligne 5 du bas de la p.~209, remplacer « localement
de type fini » par « localement de présentation finie ». L'énoncé
actuel de ({\hyperref[canon-index.ega4.c04.thm.14-4-1]{IV, 14.4.1}}) est inexact,
même lorsque $f:X\to Y$ est une immersion fermée, $X$ étant réduit à
un seul point $y$ (fermé) de $Y$, et $\sh{O}_{Y,y}$ étant un corps. Il
y a en effet des anneaux $A$ tels que $Y=\operatorname{Spec}(A)$ soit
compact, sans point isolé et totalement discontinu, tout idéal premier
$\mathfrak p$ de $A$ étant maximal et $A_\mathfrak p$ étant un corps
(Bourbaki, \emph{Alg. comm.}, chap.~II, \S4, exerc.~16 et 17); il est
clair que l'immersion fermée $\{y\}\to Y$ n'est pas un morphisme
ouvert, mais vérifie les conditions b) et b') de
({\hyperref[canon-index.ega4.c04.thm.14-4-1]{IV, 14.4.1}}).

\paragraph{(Err\textsubscript{IV}, 38).}
Dans (IV, 14.4.1.3), ligne 14 de la p.~211, remplacer « localement
quasi-fini et dominant » par « localement quasi-fini, localement de
présentation finie et dominant ».

\paragraph{(Err\textsubscript{IV}, 39).}
Dans (IV, 14.4.2), ligne 10 du bas de la p.~211, remplacer « localement
de type fini » par « localement de présentation finie ».

\paragraph{(Err\textsubscript{IV}, 40).}
Dans (IV, 14.4.3), ligne 18 de la p.~212, remplacer « localement de
type fini » par « localement de présentation finie ».

\paragraph{(Err\textsubscript{IV}, 41).}
Dans ({\hyperref[canon-index.ega4.c04.cor.14-4-4]{IV, 14.4.4}}), lignes 15 et 16 du
bas de la p.~212, remplacer « localement de type fini » par
« localement de présentation finie ». Lignes 10 et 11 du bas, remplacer
« équidimensionnel » par « localement équidimensionnel ».

\paragraph{(Err\textsubscript{IV}, 42).}
Dans (IV, 14.4.8), après la ligne 1 de la p.~214, ajouter :
\begin{quote}
\emph{Si de plus l'ensemble des composantes irréductibles de $Y$ est
localement fini, ces deux conditions sont équivalentes aux suivantes :}
\end{quote}

\paragraph{(Err\textsubscript{IV}, 43).}
Dans (IV, 14.4.10), ligne 15 de la p.~215, remplacer : quelconque, par :
dont l'ensemble des composantes irréductibles est localement fini.
Ligne 16 de la p.~215, remplacer : de type fini, par : de présentation
finie.

\paragraph{(Err\textsubscript{IV}, 44).}
Dans (IV, 15.5.6.1), ligne 15 du bas de la p.~234, après la première
phrase, ajouter :

\begin{quote}
On peut se limiter au cas où dans $X$ toute partie fermée irréductible
n'admet qu'un seul point générique. En effet, dans le cas contraire on
considère l'espace $Y=X/R$,

